% GENERATED FILE. Edit canonical lessons, then run npm run generate:books.
% canonical-source-hash: fnv1a64:4e1bea412761de3d
% canonical-lessons: MR-C326-samrat, MR-C326-simhasan, MR-C326-mukut, MR-C326-rajkanya, MR-C326-pradhan

\chapter{Emperors and Princesses}
\label{ch:mr-emperors-and-princesses}
\hlchaptermodality{326}

\begin{chapteropening}
\textbf{By the end of this chapter:} \emph{I can talk about an emperor, a throne and a crown.}

Complete the last lesson of chapter 326: I can talk about an emperor, a throne and a crown.
\end{chapteropening}

\section[\texorpdfstring{\mr{सम्राट} (samrāṭ)}{सम्राट (samrāṭ)}]{\mr{सम्राट} (samrāṭ) — an emperor}

\label{lesson:MR-C326-samrat}



\begin{quote}
Before the new one: say the Marathi for a royal court, then the Marathi for to reign.
\end{quote}

\subsection*{You'll want to know: \mr{सम्राट}}
\textbf{\mr{सम्राट}} — \emph{samrāṭ} — \textquotedblleft{}an emperor\textquotedblright{}.

\begin{quote}
\mr{तो} \mr{सम्राट} \mr{फार} \mr{शूर} \mr{आणि} \mr{दयाळू} \mr{होता}.
\end{quote}

\begin{quote}
\emph{to samrāṭ phār śūr āṇi dayāḷū hotā.}
\end{quote}

\begin{quote}
That emperor was very brave and kind.
\end{quote}

Masculine.

\mr{साम्राज्य}, an empire, is built on the same word.

\begin{cousinweb}
\textbf{\mr{सम्राट}} (\emph{samrāṭ}) is a tatsama, taken from Sanskrit \emph{samrāj}, 'a supreme ruler', in its nominative form \emph{samrāṭ} (Monier-Williams).
\end{cousinweb}

\begin{grammarlens}[title={What you've built}]
One more word for telling what happened.
\end{grammarlens}

\subsection*{Your turn}
\begin{itemize}
\raggedright
  \item \emph{Say it:} \emph{samrāṭ}
  \item \emph{Say it:} \emph{samrāṭ}, once more
  \item \emph{Say it:} \emph{to samrāṭ phār śūr āṇi dayāḷū hotā.} — That emperor was very brave and kind.
  \item \emph{Recall:} say the Marathi for a royal court, then the Marathi for to reign, then say \emph{samrāṭ} again
\end{itemize}

\begin{tcolorbox}[breakable,colback=teal!4,colframe=teal!35!black,title={Before you move on}]
What does \mr{सम्राट} mean? (\textquotedblleft{}An emperor\textquotedblright{}.) Say it once more.
\end{tcolorbox}
\section[\texorpdfstring{\mr{सिंहासन} (siṁhāsan)}{सिंहासन (siṁhāsan)}]{\mr{सिंहासन} (siṁhāsan) — a throne}

\label{lesson:MR-C326-simhasan}



\begin{quote}
Before the new one: say the Marathi for to reign, then the Marathi for an emperor.
\end{quote}

\subsection*{You'll want to know: \mr{सिंहासन}}
\textbf{\mr{सिंहासन}} — \emph{siṁhāsan} — \textquotedblleft{}a throne\textquotedblright{}.

\begin{quote}
\mr{नवा} \mr{राजा} \mr{सिंहासनावर} \mr{बसला}.
\end{quote}

\begin{quote}
\emph{navā rājā siṁhāsanāvar baslā.}
\end{quote}

\begin{quote}
The new king sat on the throne.
\end{quote}

Neuter.

\mr{सिंहासनावर} \mr{बसणे}, to take the throne.

\begin{cousinweb}
\textbf{\mr{सिंहासन}} (\emph{siṁhāsan}) is a tatsama, taken from Sanskrit \emph{siṁhāsana}, 'a lion-seat, a throne' (Monier-Williams): \emph{siṁha}, 'a lion', and \emph{āsana}, 'a seat'.
\end{cousinweb}

\begin{grammarlens}[title={What you've built}]
One more word for telling what happened.
\end{grammarlens}

\subsection*{Your turn}
\begin{itemize}
\raggedright
  \item \emph{Say it:} \emph{siṁhāsan}
  \item \emph{Say it:} \emph{siṁhāsan}, once more
  \item \emph{Say it:} \emph{navā rājā siṁhāsanāvar baslā.} — The new king sat on the throne.
  \item \emph{Recall:} say the Marathi for to reign, then the Marathi for an emperor, then say \emph{siṁhāsan} again
\end{itemize}

\begin{tcolorbox}[breakable,colback=teal!4,colframe=teal!35!black,title={Before you move on}]
What does \mr{सिंहासन} mean? (\textquotedblleft{}A throne\textquotedblright{}.) Say it once more.
\end{tcolorbox}
\section[\texorpdfstring{\mr{मुकुट} (mukuṭ)}{मुकुट (mukuṭ)}]{\mr{मुकुट} (mukuṭ) — a crown}

\label{lesson:MR-C326-mukut}



\begin{quote}
Before the new one: say the Marathi for an emperor, then the Marathi for a throne.
\end{quote}

\subsection*{You'll want to know: \mr{मुकुट}}
\textbf{\mr{मुकुट}} — \emph{mukuṭ} — \textquotedblleft{}a crown\textquotedblright{}.

\begin{quote}
\mr{राणीच्या} \mr{मुकुटात} \mr{मोती} \mr{होते}.
\end{quote}

\begin{quote}
\emph{rāṇīcyā mukuṭāt motī hote.}
\end{quote}

\begin{quote}
There were pearls in the queen's crown.
\end{quote}

Masculine.

\mr{मुकुट} \mr{घालणे}, to wear a crown.

\begin{cousinweb}
\textbf{\mr{मुकुट}} (\emph{mukuṭ}) is a tatsama, taken from Sanskrit \emph{mukuṭa}, 'a crown, a diadem' (Monier-Williams).
\end{cousinweb}

\begin{grammarlens}[title={What you've built}]
One more word for telling what happened.
\end{grammarlens}

\subsection*{Your turn}
\begin{itemize}
\raggedright
  \item \emph{Say it:} \emph{mukuṭ}
  \item \emph{Say it:} \emph{mukuṭ}, once more
  \item \emph{Say it:} \emph{rāṇīcyā mukuṭāt motī hote.} — There were pearls in the queen's crown.
  \item \emph{Recall:} say the Marathi for an emperor, then the Marathi for a throne, then say \emph{mukuṭ} again
\end{itemize}

\begin{tcolorbox}[breakable,colback=teal!4,colframe=teal!35!black,title={Before you move on}]
What does \mr{मुकुट} mean? (\textquotedblleft{}A crown\textquotedblright{}.) Say it once more.
\end{tcolorbox}
\section[\texorpdfstring{\mr{राजकन्या} (rājkanyā)}{राजकन्या (rājkanyā)}]{\mr{राजकन्या} (rājkanyā) — a princess}

\label{lesson:MR-C326-rajkanya}



\begin{quote}
Before the new one: say the Marathi for a throne, then the Marathi for a crown.
\end{quote}

\subsection*{You'll want to know: \mr{राजकन्या}}
\textbf{\mr{राजकन्या}} — \emph{rājkanyā} — \textquotedblleft{}a princess\textquotedblright{}.

\begin{quote}
\mr{राजकन्येला} \mr{घोड्यावर} \mr{बसायला} \mr{आवडायचं}.
\end{quote}

\begin{quote}
\emph{rājkanyelā ghoḍyāvar basāylā āvḍāycaṁ.}
\end{quote}

\begin{quote}
The princess used to love riding horses.
\end{quote}

Feminine.

\mr{राज}-, royal, + \mr{कन्या}, a daughter.

\begin{cousinweb}
\textbf{\mr{राजकन्या}} (\emph{rājkanyā}) is a tatsama, taken from Sanskrit \emph{rājakanyā}, 'a king's daughter' (Monier-Williams).
\end{cousinweb}

\begin{grammarlens}[title={What you've built}]
One more word for telling what happened.
\end{grammarlens}

\subsection*{Your turn}
\begin{itemize}
\raggedright
  \item \emph{Say it:} \emph{rājkanyā}
  \item \emph{Say it:} \emph{rājkanyā}, once more
  \item \emph{Say it:} \emph{rājkanyelā ghoḍyāvar basāylā āvḍāycaṁ.} — The princess used to love riding horses.
  \item \emph{Recall:} say the Marathi for a throne, then the Marathi for a crown, then say \emph{rājkanyā} again
\end{itemize}

\begin{tcolorbox}[breakable,colback=teal!4,colframe=teal!35!black,title={Before you move on}]
What does \mr{राजकन्या} mean? (\textquotedblleft{}A princess\textquotedblright{}.) Say it once more.
\end{tcolorbox}
\section[\texorpdfstring{\mr{प्रधान} (pradhān)}{प्रधान (pradhān)}]{\mr{प्रधान} (pradhān) — a minister}

\label{lesson:MR-C326-pradhan}



\begin{quote}
Before the new one: say the Marathi for a crown, then the Marathi for a princess.
\end{quote}

\subsection*{You'll want to know: \mr{प्रधान}}
\textbf{\mr{प्रधान}} — \emph{pradhān} — \textquotedblleft{}a minister\textquotedblright{}.

\begin{quote}
\mr{राजाने} \mr{हुशार} \mr{प्रधानाचा} \mr{सल्ला} \mr{घेतला}.
\end{quote}

\begin{quote}
\emph{rājāne huśār pradhānācā sallā ghetlā.}
\end{quote}

\begin{quote}
The king took the advice of his clever minister.
\end{quote}

Masculine.

Today \mr{पंतप्रधान} is the prime minister.

\begin{cousinweb}
\textbf{\mr{प्रधान}} (\emph{pradhān}) is a tatsama, taken from Sanskrit \emph{pradhāna}, 'chief, foremost; a chief minister' (Monier-Williams).
\end{cousinweb}

\begin{grammarlens}[title={What you've built}]
One more word for telling what happened.
\end{grammarlens}

\subsection*{Your turn}
\begin{itemize}
\raggedright
  \item \emph{Say it:} \emph{pradhān}
  \item \emph{Say it:} \emph{pradhān}, once more
  \item \emph{Say it:} \emph{rājāne huśār pradhānācā sallā ghetlā.} — The king took the advice of his clever minister.
  \item \emph{Recall:} say the Marathi for a crown, then the Marathi for a princess, then say \emph{pradhān} again
\end{itemize}

\begin{tcolorbox}[breakable,colback=teal!4,colframe=teal!35!black,title={Before you move on}]
What does \mr{प्रधान} mean? (\textquotedblleft{}A minister\textquotedblright{}.) Say it once more.
\end{tcolorbox}
